% !PRINTED 461
\weberpage{461}
%<<CONT
also
\[
-\gamma_2\left(\frac{-3+\omega}{2}\right) = f(\omega)^8 - \frac{256}{f(\omega)^{16}}.
\]

Setzt man also:
\weq{20}{q = e^{-\pi\sqrt{m}},}
so erh\"alt man [\S\,24, (11)]:
\[
Z = -\gamma_2\left(\frac{-3+i\sqrt{m}}{2}\right)
\]
\[
= q^{-\frac{1}{3}} \overset{\nu}{\underset{1,\infty}{\Pi}} (1+q^{2\nu-1})^8 - \frac{256\,q^{\frac{2}{3}}}{\overset{\nu}{\underset{1,\infty}{\Pi}} (1+q^{2\nu-1})^{16}}.
\]

Wir machen nun Gebrauch von der f\"ur jedes echt gebrochene positive $x$ g\"ultigen Grenzbestimmung:
\[
1+x < e^x < \frac{1}{1-x}, \quad e^{-x} > 1-x,
\]
und erhalten:
\[
1 < \Pi(1+q^{2\nu-1}) < e^{\frac{q}{1-q^2}},
\]
woraus
\[
q^{-\frac{1}{3}} - 256\,q^{\frac{2}{3}} < Z < q^{-\frac{1}{3}} e^{\frac{8q}{1-q^2}} - 256\,q^{\frac{2}{3}} e^{-\frac{16q}{1-q^2}},
\]
und indem man die obere Grenze noch vergr\"o{\ss}ert, kann man daf\"ur auch setzen:
\[
Z < q^{-\frac{1}{3}} - 256\,q^{\frac{2}{3}} + \frac{8\,q^{\frac{2}{3}}}{1-8\,q-q^2} + \frac{2^{12}\,q^{\frac{5}{3}}}{1-q^2}.
\]
F\"ur $m=11$ ist $q$ ungef\"ahr $=2^{-15}$, woraus man ersieht, da{\ss} der Unterschied beider Grenzen:
\[
\frac{8\,q^{\frac{2}{3}}}{1-8\,q-q^2} + \frac{2^{12}\,q^{\frac{5}{3}}}{1-q^2}
\]
bereits f\"ur $m=11$ und noch mehr also f\"ur die gr\"o{\ss}eren Werte von $m$ sehr klein ist. $Z$ ist also die zun\"achst \"uber
\[
q^{-\frac{1}{3}} - 256\,q^{\frac{2}{3}}
\]
gelegene ganze Zahl, und dieser Wert, f\"ur die gr\"o{\ss}eren $m$ schon $q^{-\frac{1}{3}}$, kommen einer ganzen Zahl au{\ss}erordentlich nahe. % [sic] printed singular "dieser Wert" with plural "kommen"
Daraus berechnet man diese Zahlen\footnote{Man tut gut, sich bei diesen und vielen der sp\"ater beschriebenen Rechnungen ein- f\"ur allemal den \gesperrt{Briggischen} Logarithmus
\[
\log(\pi \log e) = 0{,}134\,934\,184\,0
\]
zu merken, bei dem man \"ubrigens meist mit 7 Dezimalen ausreicht.}:
%CONT>>
